\chapter{高中导数的基本内容及要求}


\section{导数的课程教学目标}

\begin{enumerate}[(1)]
  \item 理解导数的物理意义和几何意义，并严格定义导数的概念。
  \item 了解复合函数的求导公式。
  \item 学习如何用导数求极值.
  \item 学会借助导数求闭区间上多项式函数的最值等等\cite{普通高中数学课程标准}.
\end{enumerate}


\section{导数的概念及其几何意义}


% \subsection{问题1：高台跳水运动员的速度}

% 在一次高台跳水运动中，某运动员在运动过程中的重心相对于水面的高度 $h$（单位：$\unit{m}$ ）与起跳后的时间（单位：$\unit{s}$）存在函数关系 $-h(t)=-4.9t^2+4.8t+11 $. 如何描述运动员从起跳到入水过程中运动的快慢程度？

% \begin{analysis}
%   我们经常观察跳水运动员跳水过程，可以发现，从开始到进入水中，在上升的时候，他们的动作会变得更缓慢，而在下落时会变得更快速。我们将运动的全部过程分为许多小片段，通过每个小片段里运动员的平均速度$\overline{v}$来大致描述出他们的动作状况。
% \end{analysis}

% \begin{proof}[解析]
%   \parencite{章建跃选择性必修二}
%   我们取某段时间 $\left({t}_{1}, {t}_{2}\right)$, 则 \[\overline{{v}}=\frac{h\left(t_{2}\right)-h\left(t_{1}\right)}{t_{2}-t_{1}}=-4.9\left(t_{1}+t_{2}\right)+4.8.\]

%   如果不断缩短这一时间段的长度, 即令 $t_{2}$ 无限趋于 $t_{1}$ (令 $\Delta t=t_{2}-t_{1}$, 即有 $\Delta t$ 无限趋 近于 0)，那么 $\bar{v}$ 将无限趋近于运动员在 $t_{1}$ 时刻的瞬时速度, 那么有
%   \[
%   \overline{{v}}=\lim _{\Delta t \rightarrow 0} \frac{h\left(t_{1}+\Delta t\right)-h(\Delta t)}{\Delta t}. \qedhere
%   \]
% \end{proof}


% \subsection{问题2：抛物线的切线的斜率}

% 抛物线 ${f}({x})={x}^{2}$ 在点 ${P}_{0}(1,1)$ 处的切线该如何定义?

% 与上面相似, 为了研究抛物线 ${f}({x})={x}^{2}$ 在点 ${P}_{0}(1,1)$ 处的切线, 我们在点 ${P}_{0}$ 附近任意取点 $P\left(x, x^{2}\right)$, 考察该抛物线的割线 $P_{0} P$ 的变化情况。

% \begin{figure}[htbp]
%   \centering
%   \includegraphics[width = 6cm]{1.png}
%   \caption{利用GeoGebra软件演示割线$P_{0} P$的变化情况}
%   \label{figure:利用GeoGebra软件}
% \end{figure}

% 从图~\ref{figure:利用GeoGebra软件} 可以看出, 当点 ${P}$ 无穷接近 ${P}_{0}$ 的情况下, 割线 ${P}_{0} {P}$ 是 无穷接近于某一特定 的位置, 这个确定位置的直线 ${P}_{0} {~T}$ 被称为抛物线 ${f}({x})={x}^{2}$ 在点 ${P}_{0}(1,1)$ 处的切线。

% 接下来我们要确定抛物线在点 ${P}_{0}$ 的切线 ${P}_{0} {~T}$ 的斜率 ${k}_{0}$.

% 从图~\ref{figure:利用GeoGebra软件} 和切线的定义可以看出, 抛物线 ${f}({x})={x}^{2}$ 在点 ${P}_{0}(1,1)$ 处的切线 ${P}_{0} {~T}$ 的斜率 与割线 ${P}_{0} {P}$ 的斜率之间存在内在的关系。记 $\Delta x=x-1 \quad(\Delta {x}$ 可正可负但不能为 0）,
% 则点 ${P}$ 的坐标是 $\left(1+\Delta {x},(1+\Delta x)^{2}\right)$. 于是, 割线 ${P}_{0} {P}$ 的斜率

% \[
%   k=\frac{f(x)-f(1)}{x-1}=\frac{(1+\Delta x)^{2}-1}{(1+\Delta x)-1}=\Delta x+2.
% \]

% 我们可以用割线 ${P}_{0} {P}$ 的斜率 ${k}$ 近似地表示切线 ${P}_{0} {~T}$ 的斜率 ${k}_{0}$, 并且可以通过不 断缩短横坐标间隔 $|\Delta {x}|$ 来提高近似表示的精确度。

% 当 $\Delta x$ 无穷接近 0 时, 也就是说, 无论 ${x}$ 从横轴上 1 的左边还是右边无穷接近于 1 时, 割线 $P_{0} {P}$ 的斜率 ${k}$ 都无穷接近于 2 .

% 此时 2 称作 “当 $\Delta {x}$ 无限趋近于 0 时, $k=\frac{f(1+\Delta x)-f(1)}{\Delta x}$ 的极限”, 记为 \[\lim _{\Delta x \rightarrow 0} \frac{f(1+\Delta x)-f(1)}{\Delta x}=2.\]

% 从几何角度分析, 当横坐标间隔 $|\Delta x|$ 无穷小时, 点 $P$ 将无穷接近于点 $P_{0}$, 于是 割线 $P_{0} P$ 无穷接近于点 $P_{0}$ 处的切线 $P_{0} T$ 。那么割线 $P_{0} P$ 的斜率 $k$ 无穷接近于点 $P_{0}$ 处 的切线 ${P}_{0} {~T}$ 的斜率 ${k}_{0}$, 因此, 切线 ${P}_{0} {~T}$ 的斜率 ${k}_{0}=2$.

教材”导数的概念“这一节研究了两类变化率问题：一类是物理学中的问题，涉及平均速度和瞬时速度；另一类是几何学中的问题，涉及割线斜率和切线斜率。这两类问题都采用了由平均的变化率趋向瞬时的变化率的数学思想方法. 下面我们采用这个思想方法研究更一般的问题\cite{章建跃选择性必修二}。

对于函数 ${y}={f}({x})$, 设自变量 ${x}$ 从 ${x}_{0}$ 变化到 ${x}_{0}+\Delta {x}$, 那么, 函数值 ${y}$ 就从 ${f}\left({x}_{0}\right)$ 变 化到 ${f}\left({x}_{0}+\Delta {x}\right)$. 这时, ${x}$ 的变化量为 $\Delta {x}, {y}$ 的变化量为 $\Delta y=f\left(x_{0}+\Delta x\right)-f\left(x_{0}\right)$ \cite{章建跃选择性必修二}。

我们把比值 $\frac{\Delta {y}}{\Delta {x}}$, 即 $\frac{\Delta {y}}{\Delta x}=\frac{f\left(x_{0}+\Delta x\right)-f\left(x_{0}\right)}{\Delta x}$ 叫做函数 ${y}={f}({x})$ 从 ${x}_{0}$ 到 ${x}_{0}+\Delta {x}$ 的平均变化率。
如果当 $\Delta x \rightarrow 0$ 时, 平均变化率 $\frac{\Delta {y}}{\Delta {x}}$ 无穷接近 于一个确定的值, 即 $\frac{\Delta {y}}{\Delta {x}}$ 有极限, 则称 ${y}={f}({x})$ 在 ${x}={x}_{0}$ 处可导。把这个确定的值叫做 ${y}={f}({x})$ 在 ${x}={x}_{0}$ 处的导数, 记作
\[
f^{\prime}\left({x}_{0}\right)=\lim _{\Delta x \rightarrow 0} \frac{\Delta y}{\Delta x}=\lim _{\Delta x \rightarrow 0} \frac{f\left(x_{0}+\Delta x\right)-f\left(x_{0}\right)}{\Delta x}.
\]


\subsection{导数概念教学分析}

教科书把导数概念的构建分为两个步骤：一是确定瞬时变化率的含义，二是将其一般化，并给出导数的精确定义。这个过程蕴含了逼近的数学思想和利用已知探究法。因此，在课堂上，教师要注重对学生数学思维方式的渗透与指导。

在教材中，导数的定义和对导数概念的理解，都是建立在瞬间速度这一特定的物理模型上。教学的关键是让学生理解整个过程中使用的方法以及明确瞬时速度的含义，然后是推广到一般的情形，包括数学思想、方法的迁移。老师还可以在教学中增加一些关于变化率的问题，帮助学生推断出它们的共性. 通过实例，学生能意识到导数可以用来描述一切的瞬时变化率\cite{闫伟2015高中数学}.


\subsection{导数的几何意义}

我们已经知道在 ${x}={x}_{0}$ 时, 函数 ${y}={f}$
(x) 的瞬时变化率反映了函数 ${y}={f}$
(x) 在 $x=x_{0}$ 邻近的变化, 所以 ${f}^{\prime}$ (x) 的几何意义又是怎样的呢?

\begin{figure}[htbp]
  \centering
  \includegraphics[width = 6cm]{2.png}
  \caption{导数的几何意义}
  \label{figure:导数的几何意义}
\end{figure}

图~\ref{figure:导数的几何意义} 中割线 ${P}_{0} {P}$ 的斜率比较接近于
\[
k=\frac{f(x)-f\left(x_{0}\right)}{x-x_{0}}
\]
记 $\Delta x=x-x_{0}$, 当点 ${P}$ 沿着曲线 ${y}={f}({x})$ 无穷趋近于点 ${P}_{0}$ 时, 也就是 $\Delta x \rightarrow 0$ 时, ${k}$ 无穷趋近于函数 ${y}={f}({x})$ 在 ${x}={x}_{0}$ 这里的导数。因此函数 ${y}={f}({x})$ 在 ${x}={x}_{0}$ 处的导数 ${f}^{\prime}\left({x}_{0}\right)$ 就是切线 ${P}_{0} {~T}$ 的斜率 ${k}_{0}$, 即\cite{章建跃选择性必修二}
\[
k_{0}=\lim _{\Delta x \rightarrow 0} \frac{f\left(x_{0}+\Delta x\right)-f\left(x_{0}\right)}{\Delta x}=f^{\prime}({x}) .
\]
这就是导数的几何意义。



\section{求导公式及法则}


\subsection{基本初等函数的导数公式}

\begin{enumerate}
  \item 若 $f' ({x})={c}({c}$ 为常数 $)$, 则 ${f}^{\prime}({x})=0$;
  \item 若 ${f}({x})={x}^{{c}}({c} \in \mathbb{Q}$, 且 $c \neq 0)$, 则 ${f}^{\prime}({x})={cx}^{{c}-1}$;
  \item 若 ${f}({x})=\sin {x}$, 则 ${f}^{\prime}({x})=\cos {x}$;
  \item 若 $f(x)=\cos x$, 则 $f^{\prime}(x)=-\sin x$;
  \item 若 ${f}({x})={a}^{{x}}({a}>0$, 且 ${a} \neq 1)$, 则 ${f}^{\prime}({x})=a^{x} \ln a$; 特别地, 若 ${f}({x})=e^{x}$, 则 ${f}^{\prime}({x})=e^{x}$;
  \item 若 ${f}({x})=\log _{a} x({a}>0$, 且 ${a} \neq 1)$, 则 ${f}^{\prime}({x})=\frac{1}{x \ln a}$; 特别地, 若 ${f}({x})=\ln {x}$, 则 ${f}^{\prime}({x})=\frac{1}{{x}}$ 。
\end{enumerate}


\subsection{函数的和差积商的导数公式}

一般地, 对于两个函数 ${f}({x})$ 和 ${g}({x})$ 的和 (或差) 的导数, 我们有如下法则:
\[
[f(x)\pm g(x)]^{\prime}=f^{\prime}(x)\pm g^{\prime}(x) \text { 。 }
\]
另外, 对于函数 ${f}({x})$ 和 ${g}({x})$ 的乘积 (或商) 的导数, 我们有如下法则:
\[
\begin{aligned}
&{[f(x) g(x)]^{\prime}=f^{\prime}(x) g(x)+f(x) g^{\prime}(x)} \\
&{\left[\frac{f(x)}{g(x)}\right]^{\prime}=\frac{f^{\prime}(x) g(x)-f(x) g^{\prime}(x)}{[g(x)]^{2}} \quad (g(x) \neq 0)}
\end{aligned}
\]


\subsection{复合函数的导数公式}

一般地, 对于函数 ${y}={f}({u})$ 和 $u=g({x})$ 复合而成的函数 ${y}={f}({g}({x}))$, 它的导数与函数 ${y}={f}({u})$, ${u}={g}({x})$ 的导数间的关系为 $y_{x}^{\prime}=y_{u}^{\prime} \cdot  u_{x}^{\prime}$ 。即 ${y}$ 对 ${x}$ 的导数等于 ${y}$ 对 ${u}$ 的导数与 ${u} {}$ 对 ${x}$ 的导数的乘积\cite{章建跃选择性必修二}。



\section{函数性质的导数分析}

\subsection{导数与函数的单调性}

通常, 
\begin{enumerate}[(1)]
  \item 若 ${f}({x})$ 在区间 $(a, b)$ 内大于 0 , 则曲线 ${y}={f}({x})$ 在区间 $(a, b)$ 上每一点处的切线 斜率都大于 0 , 且该曲线处于上升状态, 故 ${f}({x})$ 在区间 $(a, b)$ 上是增函数。
  \item 若 $f(x)$ 在区间 $(a, b)$ 内小于 0 , 则曲线 ${y}={f}({x})$ 在区间 $(a, b)$ 上每一点处的切线斜率都小于 0 , 曲线处 于下降状态, 故 ${f}({x})$ 在区间 $(a, b)$ 上是减函数\cite{高存明选择性必修三}。
\end{enumerate}

% 根据上述性质，可以求出函数的单调区间。


\subsection{导数与函数的极值}

设函数 $y=f(x)$ 的定义域为 $D$. 对 $x_{0} \in D$ ， 如果对 $x_{0}$ 附近的任意 $x$, 都有\cite{高存明选择性必修三}
\begin{enumerate}[(1)]
  \item $f(x) \leq f\left(x_{0}\right)$, 则称 $x_{0}$ 为函数 $f\left(x_{0}\right)$ 的一个极大值点, 且 $f(x)$ 在 $x_{0}$ 处取极大值;
  \item $f(x) \geq f\left(x_{0}\right)$, 则称 $x_{0}$ 为函数 $f(x)$ 的一个极小值点, 且 $f(x)$ 在 $x_{0}$ 处取极小值。
\end{enumerate}

显然, 极大值 点就是在其附近函数值最大的点, 极小值点就是在其附近函数值最小的点。

% 通常, 若 $x_{0}$ 是 $y=f(x)$ 的极值点, 且 $f(x)$ 在 ${x}=x_{0}$ 处可导, 则必有 $f^{\prime}\left(x_{0}\right)=0$\cite{高存明选择性必修三}. 

